\section{Equations on projective cones}
\label{sec:projective-cone-equations}

The residual construction of \cref{sec:prescribed} has a singular-space counterpart on a projective cone. A logarithmic pole at the vertex becomes a divisor on the incidence resolution, and removing it leaves a positive residual background. The radial equation below determines both the solution and its normalization scalar. Bounded uniqueness applies because the lifted density is smooth with connected dense positive locus, even though it vanishes on the exceptional section.

Let $B\subset\mathbb P^N$ be a smooth connected, linearly nondegenerate, projectively normal variety of dimension $n-1$, where $N\geq n\geq2$. Let $X\subset\mathbb P^{N+1}$ be its projective cone, with vertex $p=[0:\cdots:0:1]$, and put $\theta_0=\omega_{\mathrm{FS}}|_X$. Write the ambient homogeneous coordinates as $[Z:W]$, with $Z=(Z_0,\ldots,Z_N)$. The incidence resolution $\pi:Y\to X$ is a smooth $\mathbb P^1$-bundle $f:Y\to B$; its section $E=\pi^{-1}(p)$ is contracted, and $\pi$ is an isomorphism off $E$. These cones are reduced, irreducible, compact and normal, hence unibranch, and their vertices are singular.

The homogeneous coordinate ring of $B$ is normal by projective normality. Adjoining $W$ preserves normality, and its degree-zero projective charts are normal. Algebraic normality passes to the associated analytic space by the algebraic--analytic normality theorem. Smoothness and connectedness of $B$ give irreducibility. In the chart $W=1$, the ideal of the affine cone has no linear terms, by linear nondegeneracy. Its tangent dimension at the origin is $N+1>n$, proving that $p$ is singular. The incidence construction is the bundle of lines in $\mathcal O_B\oplus\mathcal O_B(-1)$ and has the stated contraction and isomorphism properties.

Write
\[
\begin{gathered}
\omega_B=\omega_{\mathrm{FS}}|_B,\qquad \beta=\pi^*\theta_0,\qquad \sigma=f^*\omega_B,\\
d=\int_B\omega_B^{n-1},\qquad x=\frac{\|Z\|^2}{\|Z\|^2+|W|^2}.
\end{gathered}
\]
The same symbol $x$ denotes its pullback to $Y$. Choose $\Psi\in C^\infty(\mathbb P^N,\mathbb R)$ with $\omega_{\mathrm{FS}}+\ddc\Psi>0$ and $k\in C^\infty([0,1])$ with $k>0$. Put $\psi=\Psi|_B$ and define on $Y$
\begin{equation}\label{eq:cone-prescribed-data}
\begin{gathered}
s=\frac{\mathrm{e}^\psi x}{1-x+\mathrm{e}^\psi x},\qquad \ell=\log(1-x+\mathrm{e}^\psi x),\qquad \sigma_\psi=f^*(\omega_B+\ddc\psi),\\
B_\psi=\beta+\ddc\ell,\qquad \nu=k(s)B_\psi^n,\qquad \mu=\pi_*\nu,\\
I=\int_0^1 k(r)r^{n-1}\,\mathrm{d} r,\qquad M=ndI.
\end{gathered}
\end{equation}
These data prescribe the measure before the equation is solved. The local form and fiber integral below identify its degeneracy along $E$ and its total mass $M$.

Choose a local holomorphic lift $b(z)$ of the embedding of $B$ and write $S(z)=\sum|b_j(z)|^2$, $Z=t b(z)$. On this chart
\[
\beta=\ddc\log(1+|t|^2S),\qquad \sigma=\ddc\log S.
\]
On $0<x<1$ put $\rho=\log(|t|^2S)$, $\rho_\psi=\rho+\psi$, $h(r)=\log(1+\mathrm{e}^r)$ and $\zeta_\psi=\mathrm{d}\rho_\psi\wedge\dc\rho_\psi$. Then $x=h'(\rho)$, $s=h'(\rho_\psi)$ and $\ell=h(\rho_\psi)-h(\rho)$. Differentiation gives
\begin{equation}\label{eq:cone-radial-form}
B_\psi=s\sigma_\psi+s(1-s)\zeta_\psi,\qquad \sigma_\psi^n=0,\qquad \zeta_\psi^2=0.
\end{equation}
Near $E$, a local potential is $\log(1+\mathrm{e}^\psi|t|^2S)$; at $t=0$ its Hessian has rank one in the transverse direction. At the infinity section, with $v=1/t$, a local potential is $\log(\mathrm{e}^\psi S+|v|^2)$; its tangential Hessian is $\sigma_\psi$ and its transverse coefficient is positive. These computations show that $B_\psi$ is smooth and semipositive, positive off $E$. Since $n\geq2$, its top power vanishes on $E$.

Along a fiber, $\dc\rho_\psi=\mathrm{d}(\arg t)/(2\pi)$ and $\mathrm{d}\rho_\psi=\mathrm{d} s/[s(1-s)]$. Base components disappear after wedging with $\sigma_\psi^{n-1}$, so \eqref{eq:cone-radial-form} gives
\[
f_*\nu=nI(\omega_B+\ddc\psi)^{n-1}.
\]
Smooth closed Stokes on $B$ gives $\int_B(\omega_B+\ddc\psi)^{n-1}=d$, hence $\nu(Y)=M$. Smooth measures give $E$ zero mass. Therefore $\mu(X)=M$ and $\mu(\{p\})=0$. Taking $\psi=0$ and $k=1$ gives $\int_Y\beta^n=d$.

\subsection{Logarithmic poles at the vertex}

The logarithmic pole supplies the tangential positivity missing from $\beta$ on $E$. Fix $0<a<1$ and put $\varphi=a\log x$, with value $-\infty$ at $p$. Set
\begin{equation}\label{eq:cone-logarithmic-data}
\begin{gathered}
\beta_a=(1-a)\beta+a\sigma,\qquad V=d(1-a^n),\qquad c_a=V/M,\\
P_a(s)=\left(a^n+c_an\int_0^s k(r)r^{n-1}\,\mathrm{d} r\right)^{1/n},\\
Q_a'(s)=\frac{P_a(s)-a-(1-a)s}{s(1-s)},\qquad Q_a(1)=0,\\
q_a=a\psi+(1-a)\ell+Q_a(s),\qquad A_a=\beta_a+\ddc q_a.
\end{gathered}
\end{equation}
The endpoint values in $Q_a'$ are removable. The form $\beta_a$ is strictly positive everywhere, whereas $A_a$ is semipositive and strictly positive off $E$. On that complement one may view $A_a$ as a form on $X\setminus\{p\}$.

Choose a relatively compact coordinate ball $O\subset X\setminus\{p\}$ and a smooth compactly supported endomorphism $H$ of $T^{1,0}O$ which is $A_a$-self-adjoint and trace-free. For tangent vectors $v,w$ of type $(1,0)$ at the same point of $O$, define
\begin{equation}\label{eq:cone-deformed-background}
A_H(v,\bar w)=A_a(\exp(H)v,\bar w),\qquad \eta_H=A_H-A_a,\qquad \theta_H=\theta_0+\eta_H,
\end{equation}
extending $\eta_H$ by zero. We call $H$ admissible when
\begin{equation}\label{eq:cone-admissible-positivity}
\theta_H>0,\qquad \beta_{a,H}:=\beta_a+\pi^*\eta_H>0.
\end{equation}
These are explicit open inequalities in the prescribed data. Sufficiently small $H$ satisfies them.

\Needspace{14\baselineskip}
\begin{theorem}\label{thm:cone-logarithmic-equation}\index{Projective cone!logarithmic equation}\index{Projective cone!non-descending residual}
For every admissible $H$, the ordinary class
\[
\mathcal C_{a,H}=\{w\in\PSH(X,\theta_H):\ \varphi-C\leq w\leq\varphi+C\text{ for some }C<\infty\}
\]
contains a unique pair $c>0$, $u$ with $\sup_Xu=0$ such that
\[
\mathrm{MA}^{\mathrm{ext}}_{X,\theta_H}(u)=c\mu.
\]
It is given by $c=c_a$ and
\[
u\circ\pi=\varphi\circ\pi+q_a-K,\qquad K=\sup_Y(\varphi\circ\pi+q_a).
\]
Every class member has finite extended mass. The solution has mass $V$ and is the same for every admissible $H$. For $H=0$, every class member has mass $V$.

The solution is continuous with values in the extended real line, is smooth off $p$, and tends uniformly to $-\infty$ at $p$. Its smooth residual satisfies
\[
(q_a-K)|_E=a\psi+Q_a(0)-K.
\]
If $\psi$ is nonconstant, this residual is not the pullback of any function on $X$. Arbitrarily small admissible choices of $H$ give genuinely nonclosed $\theta_H$.

For every $w\in\mathcal C_{a,H}$ and every smooth resolution $\tau:Z\to X$, including further modifications over regular points,
\begin{equation}\label{eq:cone-resolution-pushforward}
\tau_*\left\langle\bigl(\tau^*\theta_H+\ddc(w\circ\tau)\bigr)^n\right\rangle=\mathrm{MA}^{\mathrm{ext}}_{X,\theta_H}(w).
\end{equation}
The upstairs measure is finite. Thus the measure, normalized descended solution and scalar are independent of the resolution and of every common refinement.
\end{theorem}

\begin{proof}
The divisor formula is
\begin{equation}\label{eq:cone-residual-divisor}
\beta+\ddc(\varphi\circ\pi)=a[E]+\beta_a.
\end{equation}
Indeed, $\log x=\log|t|^2+\log S-\log(1+|t|^2S)$ and $\ddc\log|t|^2=[t=0]$. At $E$ the kernels of $\beta$ and $\sigma$ are complementary. Off $E$, $\beta>0$. Thus $\beta_a$ is a smooth K\"ahler form. Fiber integration also gives
\begin{equation}\label{eq:cone-logarithmic-mass}
\int_Y\beta_a^n=nd\int_0^1[a+(1-a)x]^{n-1}(1-a)\,\mathrm{d} x=d(1-a^n)=V.
\end{equation}

From \eqref{eq:cone-logarithmic-data}, $P_a(0)=a$, $P_a(1)=1$ and
\[
P_a'(s)=\frac{c_a k(s)s^{n-1}}{P_a(s)^{n-1}}.
\]
The numerator defining $Q_a'$ vanishes at both ends, proving smoothness of $Q_a$. On $0<x<1$, a local potential for the full positive current is $F_a(\rho_\psi)$, where
\[
F_a(r)=ar+(1-a)h(r)+Q_a(h'(r)).
\]
Direct differentiation gives $F_a'=P_a(s)$ and $F_a''=c_a k(s)s^n(1-s)/P_a(s)^{n-1}$. Consequently
\begin{equation}\label{eq:cone-solved-logarithmic-form}
A_a=P_a(s)\sigma_\psi+\frac{c_a k(s)s^{n-1}}{P_a(s)^{n-1}}(B_\psi-s\sigma_\psi),\qquad A_a^n=c_a\nu.
\end{equation}
The coefficients are smooth since $P_a\geq a>0$. On $E$ the restriction is $a\sigma_\psi$; at infinity the horizontal coefficient is one and the transverse coefficient relative to $B_\psi-\sigma_\psi$ is $c_a k(1)>0$. This proves all the asserted positivity properties, including at both sections.

The functions $\varphi\circ\pi$ and $\varphi\circ\pi+q_a$ are quasi-plurisubharmonic on $Y$. By \eqref{eq:cone-residual-divisor}, their exact $\beta$-currents are $a[E]+\beta_a$ and $a[E]+A_a$, respectively, so both are positive. The descent theorem \cref{thm:hermitian-space-descent} gives $\varphi\in\PSH(X,\theta_0)$ and a potential $u_0$ with $u_0\circ\pi=\varphi\circ\pi+q_a$. The bounded residual $q_a$ puts $u_0$ in the ordinary class of $\varphi$ and gives the uniform logarithmic asymptotic at $p$. Its finite supremum is $K$ in \cref{thm:cone-logarithmic-equation}.

For an admissible $H$, self-adjointness makes $\exp H$ positive and $\det\exp H=\exp(\operatorname{tr}H)=1$. Hence
\begin{equation}\label{eq:cone-determinant-preservation}
(A_a+\pi^*\eta_H)^n=A_a^n=c_a\nu,\qquad \beta_{a,H}+\ddc(q_a-K)=A_a+\pi^*\eta_H.
\end{equation}
The full function $u=u_0-K$ is unchanged. On the regular locus its exact $\theta_H$-form is the positive form in \eqref{eq:cone-determinant-preservation}; near $p$ the background and potential are unchanged. Thus $u$ is $\theta_H$-plurisubharmonic. Formula \eqref{eq:cone-residual-divisor} with the exact deformation gives $\pi^*\theta_H+\ddc(u\circ\pi)=a[E]+A_a+\pi^*\eta_H$, proving its equation after zero extension. Its mass is $c_aM=V$. This calculation uses determinant preservation, rather than a class-mass identity for a nonclosed form.

\emph{Bounded residuals and uniqueness.}

Let $w\in\mathcal C_{a,H}$. On $Y\setminus E$, its residual $t=w\circ\pi-\varphi\circ\pi$ is bounded and $\beta_{a,H}$-plurisubharmonic. By \cref{lem:qpsh-removable-analytic}, it extends with that exact background across $E$. Conversely, for any bounded $\beta_{a,H}$-plurisubharmonic $t$, the function $\varphi\circ\pi+t$ has positive current $a[E]+\beta_{a,H}+\ddc t$ in the pulled-back background. It descends by \cref{thm:hermitian-space-descent}, and its uniform bounds relative to $\varphi\circ\pi$ put the descended potential in $\mathcal C_{a,H}$. Thus bounded residuals correspond exactly to the prescribed ordinary class.

The bounded product is finite on compact $Y$ and gives $E$ zero mass, by \cref{thm:bt-input,prop:bounded-products}. Thus
\begin{equation}\label{eq:cone-competitor-pushforward}
\mathrm{MA}^{\mathrm{ext}}_{X,\theta_H}(w)=\pi_*\bigl(\beta_{a,H}+\ddc t\bigr)^n.
\end{equation}
For $H=0$, closedness makes every positive-degree exact term in the bounded expansion integrate to zero. Thus \eqref{eq:cone-logarithmic-mass} gives mass $V$ for every competitor. The nonclosed case uses only finiteness here.

If $w$ solves with scalar $c'>0$, its lifted measure is $c'\nu$, since both sides give $E$ zero mass. Both residual solutions are continuous by \cref{cor:bounded-lp-continuity}; scalar uniqueness in \cref{prop:vanishing-exponent} gives $c'=c_a$. The density $\nu/\beta_{a,H}^n$ is smooth and nonnegative with positive locus exactly $Y\setminus E$. This locus is connected and dense: it is a bundle with connected fibers $\mathbb P^1$ minus one point over connected $B$. \Cref{thm:connected-positive-uniqueness} gives $t-(q_a-K)$ constant. The normalization fixes the constant, and both full functions equal $-\infty$ at $p$. This proves uniqueness against the ordinary class.

The uniform logarithmic asymptotic proves extended-real continuity of the full solution. Its residual restriction is $a\psi+Q_a(0)-K$. If $\psi$ is nonconstant, it cannot be a pullback from $X$, since a pullback is constant on $E$. This concerns the residual; the full potential is $-\infty$ on the whole fiber and does descend.

\Needspace{8\baselineskip}
\emph{Nonclosed deformations.}

Fix $z_0\in O$ and a smooth $A_a$-unitary frame on a smaller ball. Let $K_0=\operatorname{diag}(1,-1,0,\ldots,0)$ in this frame. Choose a real compactly supported $\chi$ with $\chi(z_0)=0$ and $\mathrm{d}\chi(z_0)$ a nonzero covector in the first two complex directions, and, for $\epsilon\in\mathbb R$, put $H=\epsilon\chi K_0$. At $z_0$,
\[
\mathrm{d}\eta_H=\epsilon\,\mathrm{d}\chi\wedge\gamma,\qquad \gamma=A_a(K_0\,\cdot,\cdot).
\]
Terms differentiating the frame and $A_a$ vanish because $\chi(z_0)=0$. The restriction of $\gamma$ to these four real directions is nondegenerate, so wedge multiplication by $\gamma$ sends a nonzero one-form to a nonzero three-form. Thus $\mathrm{d}\theta_H(z_0)\ne0$ for $\epsilon\ne0$. On the fixed compact support, both strict inequalities \eqref{eq:cone-admissible-positivity} persist for sufficiently small $|\epsilon|$, and $H\to0$ smoothly. This proves existence of arbitrarily small nonclosed choices; admissibility itself does not impose nonclosedness.

For any smooth resolution $\tau:Z\to X$, \cref{thm:space-extended-measures} identifies its pushforward product with \eqref{eq:cone-competitor-pushforward}, proving \eqref{eq:cone-resolution-pushforward}. Its total mass equals the finite mass downstairs. Since the ordinary class, equation and normalization are intrinsic, uniqueness identifies the same solution and scalar on every resolution and common refinement.
\end{proof}

\subsection{The bounded equation}

Without a logarithmic pole, the residual background is $\beta$, which degenerates along $E$. The radial construction still gives a bounded solution, and closedness supplies an energy argument for its uniqueness. Put
\begin{equation}\label{eq:cone-bounded-data}
\begin{gathered}
c_b=\frac dM=\frac1{nI},\qquad R(s)=\left(c_bn\int_0^1 k(sr)r^{n-1}\,\mathrm{d} r\right)^{1/n},\\
Q'(s)=\frac{R(s)-1}{1-s},\qquad Q(1)=0.
\end{gathered}
\end{equation}
Here $R$ is smooth and positive, $R(1)=1$, and $Q$ is smooth at both endpoints.

\Needspace{12\baselineskip}
\begin{theorem}\label{thm:cone-bounded-equation}\index{Projective cone!bounded equation}
For the same cone and prescribed measure, the bounded $\theta_0$-plurisubharmonic class contains a unique normalized pair $c>0$, $u$ solving $\mathrm{MA}^{\mathrm{ext}}_{X,\theta_0}(u)=c\mu$. The scalar is $c_b$. The smooth function $\ell+Q(s)$ on $Y$ is constant on $E$ and descends to a finite continuous potential on $X$; its supremum-zero normalization is $u$.

Every bounded competitor has mass $d$. Formula \eqref{eq:cone-resolution-pushforward}, with $\theta_0$ in place of $\theta_H$, holds on every resolution and common refinement. If $\Psi=0$ and $k$ is nonconstant, or if $\psi$ is nonconstant, the solution is nonconstant.
\end{theorem}

\begin{proof}
Set $P(s)=sR(s)$ with the data in \eqref{eq:cone-bounded-data}. Then
\begin{equation}\label{eq:cone-bounded-radial-identity}
P(s)^n=c_bn\int_0^s k(r)r^{n-1}\,\mathrm{d} r,\qquad P(0)=0,\qquad P(1)=1.
\end{equation}
Smoothness and positivity of $R$ follow from its integral formula. Since $R(1)=1$, $Q$ is smooth at $s=1$, and its derivative is smooth at zero as well. Put $U=\ell+Q(s)$. It is smooth on $Y$ and equals $Q(0)$ on $E$. On $0<x<1$, a local potential for $\beta+\ddc U$ is $F(\rho_\psi)$, where $F(r)=h(r)+Q(h'(r))$. Differentiation gives $F'=P(s)$ and $F''=c_bk(s)s^n(1-s)/P(s)^{n-1}$.

\Needspace{5\baselineskip}
For $A=\beta+\ddc U$,
\begin{equation}\label{eq:cone-solved-bounded-form}
A=P(s)\sigma_\psi+\frac{c_b k(s)}{R(s)^{n-1}}(B_\psi-s\sigma_\psi),\qquad A^n=c_b\nu.
\end{equation}
The formula shows smoothness at zero directly. At infinity the horizontal coefficient is one and the transverse coefficient is positive. Thus $A$ is semipositive on $Y$ and strictly positive on $Y\setminus E$. By \cref{thm:hermitian-space-descent}, $U=u_0\circ\pi$ for $u_0\in\PSH(X,\theta_0)$. The proper surjection $\pi$ is a quotient map, so smoothness of $U$ makes $u_0$ continuous. Subtracting $K_b=\sup_Xu_0$ gives the required solution.

For any bounded $\theta_0$-plurisubharmonic $w$, the pullback $t=w\circ\pi$ is bounded $\beta$-plurisubharmonic. Its product is analytic-null and pushes forward to $\mathrm{MA}^{\mathrm{ext}}_{X,\theta_0}(w)$. Integrating the bounded closed expansion gives $\int_Y(\beta+\ddc t)^n=\int_Y\beta^n=d$, fixing the scalar of any solution to be $c_b=d/M$. The resolution identity follows from \cref{thm:space-extended-measures}. If $w$ solves, put $q=U-K_b$, $D=\beta+\ddc t$ and $v=t-q$. The bounded integration-by-parts calculus of \cref{sec:psh-local-capacity} gives
\[
0=\int_Yv(D^n-A^n)=-\sum_{j=0}^{n-1}\int_Y\mathrm{d}v\wedge\dc v\wedge D^j\wedge A^{n-1-j}.
\]
Every energy is nonnegative, so the term with the smooth factor $A^{n-1}$ vanishes. By the weak-gradient interpretation, strict positivity of $A$ on each relatively compact chart of $Y\setminus E$ forces $\mathrm{d} v=0$ there. Connectedness makes $v$ constant almost everywhere on that complement. The analytic divisor has ordinary volume zero, so \cref{lem:qpsh-almost-everywhere-order} gives equality everywhere on $Y$. Surjectivity and normalization give $w=u$ on $X$.

If $\Psi=0$ and the solution is constant, then $U=Q(x)$ is constant, so $Q'=0$ and $R=1$. Differentiating \eqref{eq:cone-bounded-radial-identity} gives $c_bk=1$, forcing $k$ constant. If $\psi$ is nonconstant, the restriction of $U$ to the infinity section is $\psi$, so the solution is nonconstant for every allowed $k$.
\end{proof}

\Needspace{12\baselineskip}
\begin{example}\label{ex:cone-angular-radial-cancellation}
Nonconstant $k$ alone does not imply nonconstant solution for unrestricted angular data. Take $\Psi=\log2$ and $k(s)=2/(2-s)^{n+1}$. Then $s=2x/(1+x)$, $\ell=\log(1+x)$ and
\[
\frac{B_\psi^n}{\beta^n}=\frac{s^n(1-s)}{x^n(1-x)}=\frac{(2-s)^{n+1}}2.
\]
Thus $\nu=\beta^n$. Since $\frac{\mathrm{d}}{\mathrm{d} s}(s/(2-s))^n=nk(s)s^{n-1}$, one has $I=1/n$, $M=d$ and $c_b=1$. The endpoint formulas give $R=1/(2-s)$ and $Q(s)=\log(2-s)$, hence $U=\log2$. The normalized solution is zero, although $k$ is nonconstant.
\end{example}

\subsection{The quadric cone}\label{sec:quadric-cone-equations}\index{Projective cone!quadric cone}

Let $X=\{Z_0Z_1=Z_2^2\}\subset\mathbb P^3$. Its base conic is the quadratic Veronese image of $\mathbb P^1$, so $n=d=2$. Its homogeneous ring $\mathbb C[\xi^2,\xi\eta,\eta^2]$ is the invariant ring of the sign involution on the normal polynomial ring, and is normal; thus the base is projectively normal. In the chart
\[
\pi(z,t)=(t,tz^2,tz),\qquad S=1+|z|^2+|z|^4,\qquad L=1+|t|^2S,
\]
one has $\beta=\ddc\log L$. Differentiation gives
\[
\det((\log L)_{j\bar k})=\frac{|t|^2(SS_{z\bar z}-|S_z|^2)}{L^3}=\frac{|t|^2(1+4|z|^2+|z|^4)}{L^3}.
\]
Writing $\lambda_4$ for Euclidean volume gives
\begin{equation}\label{eq:cone-quadric-density}
\frac{\beta^2}{2}=\frac{|t|^2(1+4|z|^2+|z|^4)}{\pi^2(1+|t|^2(1+|z|^2+|z|^4))^3}\,\mathrm{d}\lambda_4.
\end{equation}
The transverse vanishing order on $E=\{t=0\}$ is exactly two. For $\Psi=0$, $k=1/2$, one has $I=1/4$ and $M=1$, so $\mu_0=\theta_0^2/2$, extended by zero at $p$, is an ambient-volume probability. The logarithmic scalar and solution mass are $c_a=2(1-a^2)$ for every admissible deformation, including the genuinely nonclosed choices. This exact normalization follows from determinant preservation.

\Needspace{8\baselineskip}
For $\Psi=0$, $k(s)=1+s$, one has $I=5/6$, $M=10/3$ and $c_b=3/5$. Put $R(x)=\sqrt{(3+2x)/5}$ and $r_0=\sqrt{3/5}$. The normalized bounded solution is
\begin{equation}\label{eq:cone-quadric-bounded-solution}
u(x)=-2(R(x)-r_0)+2\log\frac{1+R(x)}{1+r_0}.
\end{equation}
Indeed, $u'=-2/[5(1+R)]=(R-1)/(1-x)$ with its removable value at one. It is $Q(x)-Q(0)$, strictly decreasing in $x$, with $u(p)=\sup_Xu=0$. The function of $x$ is smooth on a neighborhood of $[0,1]$, and $x$ is ambient smooth, so $u$ is a smooth ambient restriction. Its prescribed measure is $\mu_1=(1+x)\theta_0^2$, of mass $10/3$, and $\mathrm{MA}^{\mathrm{ext}}_{X,\theta_0}(u)=(3/5)\mu_1$, of mass two. Formula \eqref{eq:cone-quadric-density} also verifies the exceptional vanishing of its lift.

For a real parameter $\epsilon$, the function $\Psi=\epsilon|Z_0|^2/(|Z_0|^2+|Z_1|^2+|Z_2|^2)$ is smooth on $\mathbb P^2$ and preserves $\omega_{\mathrm{FS}}+\ddc\Psi>0$ for sufficiently small $|\epsilon|$. Its restriction to the conic takes values $\epsilon$ and zero at $[1:0:0]$ and $[0:1:0]$. For small nonzero $\epsilon$, the logarithmic residual is nonconstant on $E$ and fails to descend. Supported nonclosed deformations leave that residual, and the full descended solution, unchanged.
